\begin{lemma}[Missing colors give chain ends]
Let $\alpha\ne\beta$, let $\varphi$ be proper, and suppose that $\alpha$ is
missing at a vertex $v$.  Then the vertex-node $v$ has degree at most one in
$I_G(\varphi;\alpha,\beta)$.
\end{lemma}
\begin{proof}
By \noderef{MissingColors}, no edge incident with $v$ is colored
$\alpha$.  By the definition in \noderef{TwoColorIncidenceGraph}, every
neighbor of the vertex-node $v$ is therefore the identity of an incident
$\beta$-colored edge.  There cannot be two distinct such identities:
they would share $v$, hence be adjacent by \noderef{AdjacentEdges}, and
their common color would violate \noderef{IsProperEdgeColoring}.  Thus the
degree is zero or one, as claimed.  This is consistent with the global
maximum-degree bound \noderef{TwoColorIncidenceDegreeAtMostTwo}.
\end{proof}
